% Part: methods
% Chapter: proofs
% Section: proof-by-contradiction

\documentclass[../../../include/open-logic-section]{subfiles}

\begin{document}

\olfileid{mth}{prf}{con}

\olsection{व्याघाताने सिद्धता}

व्याघाताने सिद्धता हा प्रथमतः नकारात्मक विधाने सिद्ध करण्यासाठी वापरला जाणारा अनुमानाचा नमुना आहे. एखादे~$p$ विधान \emph{असत्य} आहे हे दाखवायचे आहे असे समजा; म्हणजे~$\lnot p$ दाखवायचे आहे. यासाठी सर्वात उपयोगी मार्ग असा: (a) $p$~सत्य आहे असे समजा आणि (b) या गृहीतकावरून असत्य असल्याचे माहीत असलेली एखादी गोष्ट मिळते हे दाखवा. “असत्य असल्याचे माहीत असलेली गोष्ट” ही $p$ शीच व्याघात करणारा निष्कर्ष असेल किंवा संपूर्ण विधानाच्या दुसऱ्या एखाद्या गृहीतकाशी व्याघात करणारा निष्कर्ष असेल. उदाहरणार्थ, “$q$ असेल, तर $\lnot p$” सिद्ध करताना $q$~सत्य आहे असे गृहीत धरून त्यावरून~$\lnot p$ सिद्ध करतो. $\lnot p$ व्याघाताने सिद्ध करायचे असेल, तर~$q$ बरोबर $p$ हेदेखील गृहीत धरायचे. $p$ वरून $\lnot q$ सिद्ध करता आले, तर~$p$ या गृहीतकावरून दुसऱ्या~$q$ गृहीतकाशी व्याघात करणारा निष्कर्ष मिळतो हे दाखवले आहे. कारण $q$~आणि $\lnot q$ दोन्ही सत्य असू शकत नाहीत. अर्थात, व्याघात सिद्ध करताना इतर अनुमानाचे नमुने वापरावे लागतील आणि व्याख्याही उलगडाव्या लागतील. एक उदाहरण पाहू.

\begin{prop}
$A \subseteq B$ आणि $B = \emptyset$ असतील, तर $A$ मध्ये कोणतेही !!{element}s नसतात.
\end{prop}

\begin{proof}
$A \subseteq B$ आणि $B = \emptyset$ असे समजा. $A$ मध्ये कोणतेही !!{element}s नाहीत हे दाखवायचे आहे.
\begin{quote}
हे सशर्त विधान असल्यामुळे पूर्वांग गृहीत धरून उत्तरांग सिद्ध करायचे आहे. उत्तरांग असे आहे: $A$ मध्ये कोणतेही !!{element}s नाहीत. हे आणखी स्पष्टपणे असे मांडता येते: असा~$x \in A$ आहे हे खरे नाही.
\end{quote}
$A$ मध्ये कोणतेही !!{element}s नसतात तेव्हाच आणि फक्त तेव्हा जेव्हा असा~$x$ आहे की $x \in A$, हे खरे नसते.
\begin{quote}
म्हणून प्रत्यक्षात~$\lnot p$ हे नकारात्मक विधान सिद्ध करायचे आहे असे ठरले: असा $x \in A$ आहे हे खरे नाही. व्याघाताने सिद्धता करण्यासाठी त्याच्याशी संबंधित~$p$ हे होकारात्मक विधान---म्हणजे असा $x \in A$ आहे---गृहीत धरून त्यावरून व्याघात सिद्ध करावा लागेल. “व्याघात मिळवण्यासाठी असे गृहीत धरा” किंवा नुसते “उलट समजा” असे लिहून~$p$ हे गृहीतक मांडले, की आपण व्याघाताने सिद्धता करत आहोत हे दर्शवता येते.
\end{quote}
उलट समजा: असा $x \in A$ आहे.
\begin{quote}
व्याघात मिळवण्यासाठी हे नवे गृहीतक वापरू. आणखी दोन गृहीतके आहेत: $A \subseteq B$ आणि $B = \emptyset$. पहिल्यावरून $x \in B$ मिळते:
\end{quote}
$A \subseteq B$ असल्यामुळे $x \in B$.
\begin{quote}
पण $B = \emptyset$ असल्यामुळे $B$ मधील प्रत्येक !!{element}---उदाहरणार्थ $x$---हा~$\emptyset$ मधील !!a{element} देखील असला पाहिजे.
\end{quote}
$B = \emptyset$ असल्यामुळे $x \in \emptyset$. हा व्याघात आहे; कारण व्याख्येनुसार $\emptyset$ मध्ये कोणतेही !!{element}s नसतात.
\begin{quote}
यानेच सिद्धता पूर्ण होते. घेतलेल्या गृहीतकांवरून आवश्यक व्याघात मिळाला आहे; म्हणजे ती सर्व गृहीतके एकाच वेळी खरी असू शकत नाहीत. पहिली दोन गृहीतके ($A \subseteq B$ आणि $B = \emptyset$) कायम ठेवलेली असल्यामुळे शेवटी आणलेले गृहीतक---असा $x \in A$ आहे---असत्य असले पाहिजे. मात्र पूर्ण तपशील द्यायचा असल्यास हे स्पष्ट लिहिता येते.
\end{quote}
म्हणून असा $x \in A$ आहे हे गृहीतक असत्य असले पाहिजे. त्यामुळे व्याघाताने सिद्धतेवरून $A$ मध्ये कोणतेही !!{element}s नाहीत.
\end{proof}

शास्त्रीय तर्कशास्त्रात प्रत्येक होकारात्मक विधान एका नकारात्मक विधानाशी सममूल्य असते: $p$ हे $\lnot\lnot p$ असते तेव्हाच आणि फक्त तेव्हाच सत्य असते. म्हणून व्याघाताने सिद्धता वापरून होकारात्मक विधानेही “अप्रत्यक्षपणे” सिद्ध करता येतात. ~$p$ सिद्ध करण्यासाठी त्याला $\lnot\lnot p$ हे नकारात्मक विधान समजा. $\lnot p$ वरून व्याघात सिद्ध करता आला, तर व्याघाताने सिद्धतेद्वारे $\lnot\lnot p$ प्रस्थापित होते आणि म्हणून~$p$ मिळते.

मागील उदाहरणात नकारात्मक विधान सिद्ध करायचे होते: $A$ मध्ये कोणतेही !!{element}s नाहीत. म्हणून व्याघाताने सिद्धतेसाठी घेतलेले गृहीतक---म्हणजे असा $x \in A$ आहे---हे होकारात्मक होते. त्यामुळे विचार करण्यासाठी तात्पुरता $x \in A$ मिळाला; त्याच्याविषयी विचार पुढे नेत $x \in \emptyset$ पर्यंत पोहोचलो.

होकारात्मक विधान अप्रत्यक्षपणे सिद्ध करताना व्याघात मिळवण्यासाठी नकारात्मक गृहीतक घ्यावे लागते. पण अनेकदा होकारात्मक विधान नकारात्मक रूपात आणि नकारात्मक विधान होकारात्मक रूपात सहज पुन्हा मांडता येते. मागच्या सिद्धतेत “$A$ मध्ये कोणतेही !!{element}s नाहीत” या नकारात्मक उत्तरांगाऐवजी “$A = \emptyset$” सिद्ध केले असते, तरी सिद्धता मूलतः तशीच असती. $=$ च्या व्याख्येनुसार “$A = \emptyset$” हे सामान्य विधान आहे; कारण ते उलगडल्यावर “~$A$ मधील प्रत्येक !!{element} हा~$\emptyset$ मधील !!{element} आहे आणि उलटही तसेच आहे” असे मिळते. पण ते “असा $x \in A$ आहे हे खरे नाही” या नकारात्मक विधानाशी सममूल्य आहे हे सहज दिसते.

म्हणून कधी~$p$ प्रत्यक्ष सिद्ध करण्यापेक्षा $\lnot p$ हे गृहीतक घेऊन काम करणे सोपे असते. पुढील उदाहरणांप्रमाणे प्रत्यक्ष सिद्धता तितकीच सोपी किंवा अधिक सोपी असली, तरी काही जण अप्रत्यक्ष मार्ग पसंत करतात. द्विनकरणामुळे गोंधळ होत असेल, तर व्याघाताने सिद्धता म्हणजे \emph{विरुद्ध} विधान गृहीत धरून व्याघात सिद्ध करणे, असा विचार करा. म्हणून $\lnot p$ व्याघाताने सिद्ध करणे म्हणजे~$p$ गृहीत धरून व्याघात सिद्ध करणे; आणि~$p$ व्याघाताने सिद्ध करणे म्हणजे~$\lnot p$ वरून व्याघात सिद्ध करणे.

\begin{prop}
$A \subseteq A \cup B$.
\end{prop}

\begin{proof}
$A \subseteq A \cup B$ दाखवायचे आहे.
\begin{quote}
प्रथमदर्शनी हे होकारात्मक विधान आहे: प्रत्येक $x \in A$ हा $A \cup B$ मध्येही असतो. त्याचे नकरण असे: काही $x \in A$ हा $\notin A \cup B$ असतो. म्हणून हे नकरण गृहीत धरून त्यावरून व्याघात मिळतो असे दाखवल्यास विधान अप्रत्यक्षपणे सिद्ध होते.
\end{quote}
उलट समजा; म्हणजे $A \nsubseteq A \cup B$.
\begin{quote}
$A \subseteq A \cup B$ ची व्याख्या अशी आहे: प्रत्येक $x \in A$ हा $\in A \cup B$ देखील असतो. $A \nsubseteq A \cup B$ चा अर्थ समजण्यासाठी उलगडलेल्या व्याख्येवर काही मूलभूत तार्किक क्रिया कराव्या लागतात: प्रत्येक $x \in A$ हा $\in A \cup B$ देखील असतो हे असत्य असते तेव्हाच आणि फक्त तेव्हा जेव्हा असा \emph{काही}~$x \in A$ असतो की तो $\notin A \cup B$ असतो. येथे फार काळजी घ्यावी: “सर्व $A$ हे $B$ आहेत” अशा विधानाचे नकरण चुकीचे उलगडल्यामुळे अनेक विद्यार्थ्यांच्या व्याघाताने सिद्धता करण्याच्या प्रयत्नांत चूक होते. म्हणजे $A \nsubseteq A \cup B$ हे तेव्हाच आणि फक्त तेव्हा सत्य असते जेव्हा असा $x$ आहे की $x \in A$ आणि $x \notin A \cup B$. यापुढे सिद्धता सोपी आहे.
\end{quote}
म्हणून असा $x \in A$ आहे की $x \notin A \cup B$. $\cup$ च्या व्याख्येनुसार $x \in A \cup B$ हे $x \in A$ किंवा $x \in B$ असते तेव्हाच आणि फक्त तेव्हाच सत्य असते. $x \in A$ असल्यामुळे $x \in A \cup B$ मिळते. हा $x \notin A \cup B$ या गृहीतकाशी व्याघात आहे.
\end{proof}

\begin{prob}
$A \cap B \subseteq A$ हे \emph{अप्रत्यक्षपणे} सिद्ध करा.
\end{prob}

\begin{prop}
$A \subseteq B$ आणि $B \subseteq C$ असतील, तर $A \subseteq C$.
\end{prop}

\begin{proof}
$A \subseteq B$ आणि $B \subseteq C$ असे समजा. $A \subseteq C$ दाखवायचे आहे.
\begin{quote}
अप्रत्यक्ष मार्ग वापरू: जे प्रस्थापित करायचे आहे त्याचे नकरण गृहीत धरू.
\end{quote}
उलट समजा; म्हणजे $A \nsubseteq C$.
\begin{quote}
आधीप्रमाणे, $A \nsubseteq C$ हे तेव्हाच आणि फक्त तेव्हा सत्य असते जेव्हा प्रत्येक $x \in A$ हा $\in C$ देखील असतो हे खरे नसते; म्हणजे काही $x \in A$ हा $\notin C$ असतो. सराव झाल्यावर अशी नकरणे उलगडण्यासाठी फार विचार करावा लागणार नाही.
\end{quote}
म्हणजे असा~$x$ आहे की $x \in A$ आणि $x \notin C$.
\begin{quote}
आता यावरून व्याघात मिळवता येईल. अर्थात, त्यासाठी इतर दोन गृहीतके वापरावी लागतील.
\end{quote}
$A \subseteq B$ असल्यामुळे $x \in B$. $B \subseteq C$ असल्यामुळे $x \in C$. पण हा $x \notin C$ शी व्याघात आहे.
\end{proof}

\begin{prop}
$A \cup B = A \cap B$ असेल, तर $A = B$.
\end{prop}

\begin{proof}
$A \cup B = A \cap B$ असे समजा. $A = B$ दाखवायचे आहे.
\begin{quote}
आता सुरुवातीची पद्धत परिचित आहे:
\end{quote}
व्याघात मिळवण्यासाठी $A \neq B$ असे गृहीत धरा.
\begin{quote}
व्याघाताने सिद्धतेसाठीचे गृहीतक $A \neq B$ आहे. $A = B$ हे $A \subseteq B$ आणि $B \subseteq A$ असतात तेव्हाच आणि फक्त तेव्हाच सत्य असल्यामुळे, $A \neq B$ हे $A \nsubseteq B$ \emph{किंवा} $B \nsubseteq A$ असते तेव्हाच आणि फक्त तेव्हाच सत्य असते. नकरणांवरील क्रिया काळजीपूर्वक करणे किती महत्त्वाचे आहे ते पाहा!{} या विकल्पयोगावरून व्याघात सिद्ध करण्यासाठी प्रकरणांनुसार सिद्धता वापरून प्रत्येक प्रकरणात व्याघात मिळतो हे दाखवू.
\end{quote}
$A \neq B$ हे $A \nsubseteq B$ किंवा $B \nsubseteq A$ असते तेव्हाच आणि फक्त तेव्हाच सत्य असते. प्रकरणे वेगळी पाहू.
\begin{quote}
पहिल्या प्रकरणात $A \nsubseteq B$ गृहीत धरतो; म्हणजे काही $x$ साठी $x \in A$, पण $\notin B$. $A \cap B$ ची व्याख्या $A$ आणि $B$ यांच्या समान !!{element}s चा संच अशी आहे. म्हणून एखादी वस्तू त्यांपैकी एका संचात नसेल, तर ती छेदात नाही. $A \cup B$ हा $A$ बरोबर $B$ घेऊन मिळालेला संच आहे. म्हणून कोणत्याही एका संचातील वस्तू संयोगातही असते. यावरून $x \in A \cup B$, पण $x \notin A \cap B$, असे मिळते; म्हणून $A \cap B \neq A \cup B$.
\end{quote}

प्रकरण 1: $A \nsubseteq B$. मग काही $x$ साठी $x \in A$, पण $x \notin B$. $x \notin B$ असल्यामुळे $x \notin A \cap B$. $x \in A$ असल्यामुळे $x \in A \cup B$. म्हणून $A \cap B \neq A \cup B$; हा $A \cap B = A \cup B$ या गृहीतकाशी व्याघात आहे.

प्रकरण 2: $B \nsubseteq A$. मग काही $y$ साठी $y \in B$, पण $y \notin A$. आधीप्रमाणे $y \in A \cup B$, पण $y \notin A \cap B$. म्हणून $A \cap B \neq A \cup B$; हा पुन्हा $A \cap B = A \cup B$ शी व्याघात आहे.
\end{proof}

\end{document}
